% Options for packages loaded elsewhere
\PassOptionsToPackage{unicode}{hyperref}
\PassOptionsToPackage{hyphens}{url}
\documentclass[
]{article}
\usepackage{xcolor}
\usepackage[margin=1in]{geometry}
\usepackage{amsmath,amssymb}
\setcounter{secnumdepth}{-\maxdimen} % remove section numbering
\usepackage{iftex}
\ifPDFTeX
  \usepackage[T1]{fontenc}
  \usepackage[utf8]{inputenc}
  \usepackage{textcomp} % provide euro and other symbols
\else % if luatex or xetex
  \usepackage{unicode-math} % this also loads fontspec
  \defaultfontfeatures{Scale=MatchLowercase}
  \defaultfontfeatures[\rmfamily]{Ligatures=TeX,Scale=1}
\fi
\usepackage{lmodern}
\ifPDFTeX\else
  % xetex/luatex font selection
\fi
% Use upquote if available, for straight quotes in verbatim environments
\IfFileExists{upquote.sty}{\usepackage{upquote}}{}
\IfFileExists{microtype.sty}{% use microtype if available
  \usepackage[]{microtype}
  \UseMicrotypeSet[protrusion]{basicmath} % disable protrusion for tt fonts
}{}
\makeatletter
\@ifundefined{KOMAClassName}{% if non-KOMA class
  \IfFileExists{parskip.sty}{%
    \usepackage{parskip}
  }{% else
    \setlength{\parindent}{0pt}
    \setlength{\parskip}{6pt plus 2pt minus 1pt}}
}{% if KOMA class
  \KOMAoptions{parskip=half}}
\makeatother
\usepackage{color}
\usepackage{fancyvrb}
\newcommand{\VerbBar}{|}
\newcommand{\VERB}{\Verb[commandchars=\\\{\}]}
\DefineVerbatimEnvironment{Highlighting}{Verbatim}{commandchars=\\\{\}}
% Add ',fontsize=\small' for more characters per line
\usepackage{framed}
\definecolor{shadecolor}{RGB}{248,248,248}
\newenvironment{Shaded}{\begin{snugshade}}{\end{snugshade}}
\newcommand{\AlertTok}[1]{\textcolor[rgb]{0.94,0.16,0.16}{#1}}
\newcommand{\AnnotationTok}[1]{\textcolor[rgb]{0.56,0.35,0.01}{\textbf{\textit{#1}}}}
\newcommand{\AttributeTok}[1]{\textcolor[rgb]{0.13,0.29,0.53}{#1}}
\newcommand{\BaseNTok}[1]{\textcolor[rgb]{0.00,0.00,0.81}{#1}}
\newcommand{\BuiltInTok}[1]{#1}
\newcommand{\CharTok}[1]{\textcolor[rgb]{0.31,0.60,0.02}{#1}}
\newcommand{\CommentTok}[1]{\textcolor[rgb]{0.56,0.35,0.01}{\textit{#1}}}
\newcommand{\CommentVarTok}[1]{\textcolor[rgb]{0.56,0.35,0.01}{\textbf{\textit{#1}}}}
\newcommand{\ConstantTok}[1]{\textcolor[rgb]{0.56,0.35,0.01}{#1}}
\newcommand{\ControlFlowTok}[1]{\textcolor[rgb]{0.13,0.29,0.53}{\textbf{#1}}}
\newcommand{\DataTypeTok}[1]{\textcolor[rgb]{0.13,0.29,0.53}{#1}}
\newcommand{\DecValTok}[1]{\textcolor[rgb]{0.00,0.00,0.81}{#1}}
\newcommand{\DocumentationTok}[1]{\textcolor[rgb]{0.56,0.35,0.01}{\textbf{\textit{#1}}}}
\newcommand{\ErrorTok}[1]{\textcolor[rgb]{0.64,0.00,0.00}{\textbf{#1}}}
\newcommand{\ExtensionTok}[1]{#1}
\newcommand{\FloatTok}[1]{\textcolor[rgb]{0.00,0.00,0.81}{#1}}
\newcommand{\FunctionTok}[1]{\textcolor[rgb]{0.13,0.29,0.53}{\textbf{#1}}}
\newcommand{\ImportTok}[1]{#1}
\newcommand{\InformationTok}[1]{\textcolor[rgb]{0.56,0.35,0.01}{\textbf{\textit{#1}}}}
\newcommand{\KeywordTok}[1]{\textcolor[rgb]{0.13,0.29,0.53}{\textbf{#1}}}
\newcommand{\NormalTok}[1]{#1}
\newcommand{\OperatorTok}[1]{\textcolor[rgb]{0.81,0.36,0.00}{\textbf{#1}}}
\newcommand{\OtherTok}[1]{\textcolor[rgb]{0.56,0.35,0.01}{#1}}
\newcommand{\PreprocessorTok}[1]{\textcolor[rgb]{0.56,0.35,0.01}{\textit{#1}}}
\newcommand{\RegionMarkerTok}[1]{#1}
\newcommand{\SpecialCharTok}[1]{\textcolor[rgb]{0.81,0.36,0.00}{\textbf{#1}}}
\newcommand{\SpecialStringTok}[1]{\textcolor[rgb]{0.31,0.60,0.02}{#1}}
\newcommand{\StringTok}[1]{\textcolor[rgb]{0.31,0.60,0.02}{#1}}
\newcommand{\VariableTok}[1]{\textcolor[rgb]{0.00,0.00,0.00}{#1}}
\newcommand{\VerbatimStringTok}[1]{\textcolor[rgb]{0.31,0.60,0.02}{#1}}
\newcommand{\WarningTok}[1]{\textcolor[rgb]{0.56,0.35,0.01}{\textbf{\textit{#1}}}}
\usepackage{graphicx}
\makeatletter
\newsavebox\pandoc@box
\newcommand*\pandocbounded[1]{% scales image to fit in text height/width
  \sbox\pandoc@box{#1}%
  \Gscale@div\@tempa{\textheight}{\dimexpr\ht\pandoc@box+\dp\pandoc@box\relax}%
  \Gscale@div\@tempb{\linewidth}{\wd\pandoc@box}%
  \ifdim\@tempb\p@<\@tempa\p@\let\@tempa\@tempb\fi% select the smaller of both
  \ifdim\@tempa\p@<\p@\scalebox{\@tempa}{\usebox\pandoc@box}%
  \else\usebox{\pandoc@box}%
  \fi%
}
% Set default figure placement to htbp
\def\fps@figure{htbp}
\makeatother
\setlength{\emergencystretch}{3em} % prevent overfull lines
\providecommand{\tightlist}{%
  \setlength{\itemsep}{0pt}\setlength{\parskip}{0pt}}
\usepackage{bookmark}
\IfFileExists{xurl.sty}{\usepackage{xurl}}{} % add URL line breaks if available
\urlstyle{same}
\hypersetup{
  pdftitle={Foundations for statistical inference - Sampling distributions},
  pdfauthor={Daanish Rasheed},
  hidelinks,
  pdfcreator={LaTeX via pandoc}}

\title{Foundations for statistical inference - Sampling distributions}
\author{Daanish Rasheed}
\date{}

\begin{document}
\maketitle

In this lab, you will investigate the ways in which the statistics from
a random sample of data can serve as point estimates for population
parameters. We're interested in formulating a \emph{sampling
distribution} of our estimate in order to learn about the properties of
the estimate, such as its distribution.

\protect\phantomsection\label{boxedtext}
\textbf{Setting a seed:} We will take some random samples and build
sampling distributions in this lab, which means you should set a seed at
the start of your lab. If this concept is new to you, review the lab on
probability.

\subsection{Getting Started}\label{getting-started}

\subsubsection{Load packages}\label{load-packages}

In this lab, we will explore and visualize the data using the
\textbf{tidyverse} suite of packages. We will also use the
\textbf{infer} package for resampling.

Let's load the packages.

\begin{Shaded}
\begin{Highlighting}[]
\FunctionTok{library}\NormalTok{(tidyverse)}
\FunctionTok{library}\NormalTok{(openintro)}
\FunctionTok{library}\NormalTok{(infer)}
\end{Highlighting}
\end{Shaded}

\subsubsection{The data}\label{the-data}

A 2019 Gallup report states the following:

\begin{quote}
The premise that scientific progress benefits people has been embodied
in discoveries throughout the ages -- from the development of
vaccinations to the explosion of technology in the past few decades,
resulting in billions of supercomputers now resting in the hands and
pockets of people worldwide. Still, not everyone around the world feels
science benefits them personally.

\textbf{Source:}
\href{https://news.gallup.com/opinion/gallup/268121/world-science-day-knowledge-power.aspx}{World
Science Day: Is Knowledge Power?}
\end{quote}

The Wellcome Global Monitor finds that 20\% of people globally do not
believe that the work scientists do benefits people like them. In this
lab, you will assume this 20\% is a true population proportion and learn
about how sample proportions can vary from sample to sample by taking
smaller samples from the population. We will first create our population
assuming a population size of 100,000. This means 20,000 (20\%) of the
population think the work scientists do does not benefit them personally
and the remaining 80,000 think it does.

\begin{Shaded}
\begin{Highlighting}[]
\NormalTok{global\_monitor }\OtherTok{\textless{}{-}} \FunctionTok{tibble}\NormalTok{(}
  \AttributeTok{scientist\_work =} \FunctionTok{c}\NormalTok{(}\FunctionTok{rep}\NormalTok{(}\StringTok{"Benefits"}\NormalTok{, }\DecValTok{80000}\NormalTok{), }\FunctionTok{rep}\NormalTok{(}\StringTok{"Doesn\textquotesingle{}t benefit"}\NormalTok{, }\DecValTok{20000}\NormalTok{))}
\NormalTok{)}
\end{Highlighting}
\end{Shaded}

The name of the data frame is \texttt{global\_monitor} and the name of
the variable that contains responses to the question \emph{``Do you
believe that the work scientists do benefit people like you?''} is
\texttt{scientist\_work}.

We can quickly visualize the distribution of these responses using a bar
plot.

\begin{Shaded}
\begin{Highlighting}[]
\FunctionTok{ggplot}\NormalTok{(global\_monitor, }\FunctionTok{aes}\NormalTok{(}\AttributeTok{x =}\NormalTok{ scientist\_work)) }\SpecialCharTok{+}
  \FunctionTok{geom\_bar}\NormalTok{() }\SpecialCharTok{+}
  \FunctionTok{labs}\NormalTok{(}
    \AttributeTok{x =} \StringTok{""}\NormalTok{, }\AttributeTok{y =} \StringTok{""}\NormalTok{,}
    \AttributeTok{title =} \StringTok{"Do you believe that the work scientists do benefit people like you?"}
\NormalTok{  ) }\SpecialCharTok{+}
  \FunctionTok{coord\_flip}\NormalTok{() }
\end{Highlighting}
\end{Shaded}

\pandocbounded{\includegraphics[keepaspectratio]{Lab5a_sampling_distributions_files/figure-latex/bar-plot-pop-1.pdf}}

We can also obtain summary statistics to confirm we constructed the data
frame correctly.

\begin{Shaded}
\begin{Highlighting}[]
\NormalTok{global\_monitor }\SpecialCharTok{\%\textgreater{}\%}
  \FunctionTok{count}\NormalTok{(scientist\_work) }\SpecialCharTok{\%\textgreater{}\%}
  \FunctionTok{mutate}\NormalTok{(}\AttributeTok{p =}\NormalTok{ n }\SpecialCharTok{/}\FunctionTok{sum}\NormalTok{(n))}
\end{Highlighting}
\end{Shaded}

\begin{verbatim}
## # A tibble: 2 x 3
##   scientist_work      n     p
##   <chr>           <int> <dbl>
## 1 Benefits        80000   0.8
## 2 Doesn't benefit 20000   0.2
\end{verbatim}

\subsection{The unknown sampling
distribution}\label{the-unknown-sampling-distribution}

In this lab, you have access to the entire population, but this is
rarely the case in real life. Gathering information on an entire
population is often extremely costly or impossible. Because of this, we
often take a sample of the population and use that to understand the
properties of the population.

If you are interested in estimating the proportion of people who don't
think the work scientists do benefits them, you can use the
\texttt{sample\_n} command to survey the population.

\begin{Shaded}
\begin{Highlighting}[]
\NormalTok{samp1 }\OtherTok{\textless{}{-}}\NormalTok{ global\_monitor }\SpecialCharTok{\%\textgreater{}\%}
  \FunctionTok{sample\_n}\NormalTok{(}\DecValTok{50}\NormalTok{)}
\end{Highlighting}
\end{Shaded}

This command collects a simple random sample of size 50 from the
\texttt{global\_monitor} dataset, and assigns the result to
\texttt{samp1}. This is similar to randomly drawing names from a hat
that contains the names of all in the population. Working with these 50
names is considerably simpler than working with all 100,000 people in
the population.

\begin{enumerate}
\def\labelenumi{\arabic{enumi}.}
\tightlist
\item
  Describe the distribution of responses in this sample. How does it
  compare to the distribution of responses in the population.
  \textbf{Hint:} Although the \texttt{sample\_n} function takes a random
  sample of observations (i.e.~rows) from the dataset, you can still
  refer to the variables in the dataset with the same names. Code you
  presented earlier for visualizing and summarizing the population data
  will still be useful for the sample, however be careful to not label
  your proportion \texttt{p} since you're now calculating a sample
  statistic, not a population parameters. You can customize the label of
  the statistics to indicate that it comes from the sample.
\end{enumerate}

\textbf{Insert your answer here}

\begin{Shaded}
\begin{Highlighting}[]
\NormalTok{samp1 }\SpecialCharTok{\%\textgreater{}\%}
  \FunctionTok{count}\NormalTok{(scientist\_work) }\SpecialCharTok{\%\textgreater{}\%}
  \FunctionTok{mutate}\NormalTok{(}\AttributeTok{s =}\NormalTok{ n }\SpecialCharTok{/}\FunctionTok{sum}\NormalTok{(n))}
\end{Highlighting}
\end{Shaded}

\begin{verbatim}
## # A tibble: 2 x 3
##   scientist_work      n     s
##   <chr>           <int> <dbl>
## 1 Benefits           37  0.74
## 2 Doesn't benefit    13  0.26
\end{verbatim}

According to the results of the above code, 74\% of the sample believe
that work scientists do provide benefit to people. The sample
proportions are not exactly equal to the population proportions. The gap
is due to sampling variability.

If you're interested in estimating the proportion of all people who do
not believe that the work scientists do benefits them, but you do not
have access to the population data, your best single guess is the sample
mean.

\begin{Shaded}
\begin{Highlighting}[]
\NormalTok{samp1 }\SpecialCharTok{\%\textgreater{}\%}
  \FunctionTok{count}\NormalTok{(scientist\_work) }\SpecialCharTok{\%\textgreater{}\%}
  \FunctionTok{mutate}\NormalTok{(}\AttributeTok{p\_hat =}\NormalTok{ n }\SpecialCharTok{/}\FunctionTok{sum}\NormalTok{(n))}
\end{Highlighting}
\end{Shaded}

\begin{verbatim}
## # A tibble: 2 x 3
##   scientist_work      n p_hat
##   <chr>           <int> <dbl>
## 1 Benefits           37  0.74
## 2 Doesn't benefit    13  0.26
\end{verbatim}

Depending on which 50 people you selected, your estimate could be a bit
above or a bit below the true population proportion of 0.26. In general,
though, the sample proportion turns out to be a pretty good estimate of
the true population proportion, and you were able to get it by sampling
less than 1\% of the population.

\begin{enumerate}
\def\labelenumi{\arabic{enumi}.}
\setcounter{enumi}{1}
\tightlist
\item
  Would you expect the sample proportion to match the sample proportion
  of another student's sample? Why, or why not? If the answer is no,
  would you expect the proportions to be somewhat different or very
  different? Ask a student team to confirm your answer.
\end{enumerate}

No, because every student takes a unique random sample of 50 people,
which means that their sample proportions are not going to be the same.
However, the sample proportions will be quite similar since they all
come from the same population.

\begin{enumerate}
\def\labelenumi{\arabic{enumi}.}
\setcounter{enumi}{2}
\tightlist
\item
  Take a second sample, also of size 50, and call it \texttt{samp2}. How
  does the sample proportion of \texttt{samp2} compare with that of
  \texttt{samp1}? Suppose we took two more samples, one of size 100 and
  one of size 1000. Which would you think would provide a more accurate
  estimate of the population proportion?
\end{enumerate}

\textbf{Insert your answer here}

\begin{Shaded}
\begin{Highlighting}[]
\NormalTok{samp2 }\OtherTok{\textless{}{-}}\NormalTok{ global\_monitor }\SpecialCharTok{\%\textgreater{}\%}
  \FunctionTok{sample\_n}\NormalTok{(}\DecValTok{50}\NormalTok{)}

\NormalTok{samp2 }\SpecialCharTok{\%\textgreater{}\%}
  \FunctionTok{count}\NormalTok{(scientist\_work) }\SpecialCharTok{\%\textgreater{}\%}
  \FunctionTok{mutate}\NormalTok{(}\AttributeTok{p\_hat =}\NormalTok{ n }\SpecialCharTok{/} \FunctionTok{sum}\NormalTok{(n))}
\end{Highlighting}
\end{Shaded}

\begin{verbatim}
## # A tibble: 2 x 3
##   scientist_work      n p_hat
##   <chr>           <int> <dbl>
## 1 Benefits           42  0.84
## 2 Doesn't benefit     8  0.16
\end{verbatim}

\begin{Shaded}
\begin{Highlighting}[]
\NormalTok{samp3 }\OtherTok{\textless{}{-}}\NormalTok{ global\_monitor }\SpecialCharTok{\%\textgreater{}\%}
  \FunctionTok{sample\_n}\NormalTok{(}\DecValTok{100}\NormalTok{)}

\NormalTok{samp3 }\SpecialCharTok{\%\textgreater{}\%}
  \FunctionTok{count}\NormalTok{(scientist\_work) }\SpecialCharTok{\%\textgreater{}\%}
  \FunctionTok{mutate}\NormalTok{(}\AttributeTok{p\_hat =}\NormalTok{ n }\SpecialCharTok{/} \FunctionTok{sum}\NormalTok{(n))}
\end{Highlighting}
\end{Shaded}

\begin{verbatim}
## # A tibble: 2 x 3
##   scientist_work      n p_hat
##   <chr>           <int> <dbl>
## 1 Benefits           81  0.81
## 2 Doesn't benefit    19  0.19
\end{verbatim}

\begin{Shaded}
\begin{Highlighting}[]
\NormalTok{samp4 }\OtherTok{\textless{}{-}}\NormalTok{ global\_monitor }\SpecialCharTok{\%\textgreater{}\%}
  \FunctionTok{sample\_n}\NormalTok{(}\DecValTok{1000}\NormalTok{)}

\NormalTok{samp4 }\SpecialCharTok{\%\textgreater{}\%}
  \FunctionTok{count}\NormalTok{(scientist\_work) }\SpecialCharTok{\%\textgreater{}\%}
  \FunctionTok{mutate}\NormalTok{(}\AttributeTok{p\_hat =}\NormalTok{ n }\SpecialCharTok{/} \FunctionTok{sum}\NormalTok{(n))}
\end{Highlighting}
\end{Shaded}

\begin{verbatim}
## # A tibble: 2 x 3
##   scientist_work      n p_hat
##   <chr>           <int> <dbl>
## 1 Benefits          798 0.798
## 2 Doesn't benefit   202 0.202
\end{verbatim}

It is expected that the sample of 1000 gives the most accurate estimate,
followed by the sample of 100. The standard error shrinks with the
square root of the sample size (about 0.057 for n = 50, 0.040 for n =
100 and 0.013 for n = 1000), so larger samples land closer to 0.20 more
often. Any single sample can still miss by chance.

Not surprisingly, every time you take another random sample, you might
get a different sample proportion. It's useful to get a sense of just
how much variability you should expect when estimating the population
mean this way. The distribution of sample proportions, called the
\emph{sampling distribution (of the proportion)}, can help you
understand this variability. In this lab, because you have access to the
population, you can build up the sampling distribution for the sample
proportion by repeating the above steps many times. Here, we use R to
take 15,000 different samples of size 50 from the population, calculate
the proportion of responses in each sample, filter for only the
\emph{Doesn't benefit} responses, and store each result in a vector
called \texttt{sample\_props50}. Note that we specify that
\texttt{replace\ =\ TRUE} since sampling distributions are constructed
by sampling with replacement.

\begin{Shaded}
\begin{Highlighting}[]
\NormalTok{sample\_props50 }\OtherTok{\textless{}{-}}\NormalTok{ global\_monitor }\SpecialCharTok{\%\textgreater{}\%}
                    \FunctionTok{rep\_sample\_n}\NormalTok{(}\AttributeTok{size =} \DecValTok{50}\NormalTok{, }\AttributeTok{reps =} \DecValTok{15000}\NormalTok{, }\AttributeTok{replace =} \ConstantTok{TRUE}\NormalTok{) }\SpecialCharTok{\%\textgreater{}\%}
                    \FunctionTok{count}\NormalTok{(scientist\_work) }\SpecialCharTok{\%\textgreater{}\%}
                    \FunctionTok{mutate}\NormalTok{(}\AttributeTok{p\_hat =}\NormalTok{ n }\SpecialCharTok{/}\FunctionTok{sum}\NormalTok{(n)) }\SpecialCharTok{\%\textgreater{}\%}
                    \FunctionTok{filter}\NormalTok{(scientist\_work }\SpecialCharTok{==} \StringTok{"Doesn\textquotesingle{}t benefit"}\NormalTok{)}
\end{Highlighting}
\end{Shaded}

And we can visualize the distribution of these proportions with a
histogram.

\begin{Shaded}
\begin{Highlighting}[]
\FunctionTok{ggplot}\NormalTok{(}\AttributeTok{data =}\NormalTok{ sample\_props50, }\FunctionTok{aes}\NormalTok{(}\AttributeTok{x =}\NormalTok{ p\_hat)) }\SpecialCharTok{+}
  \FunctionTok{geom\_histogram}\NormalTok{(}\AttributeTok{binwidth =} \FloatTok{0.02}\NormalTok{) }\SpecialCharTok{+}
  \FunctionTok{labs}\NormalTok{(}
    \AttributeTok{x =} \StringTok{"p\_hat (Doesn\textquotesingle{}t benefit)"}\NormalTok{,}
    \AttributeTok{title =} \StringTok{"Sampling distribution of p\_hat"}\NormalTok{,}
    \AttributeTok{subtitle =} \StringTok{"Sample size = 50, Number of samples = 15000"}
\NormalTok{  )}
\end{Highlighting}
\end{Shaded}

Next, you will review how this set of code works.

\begin{enumerate}
\def\labelenumi{\arabic{enumi}.}
\setcounter{enumi}{3}
\tightlist
\item
  How many elements are there in \texttt{sample\_props50}? Describe the
  sampling distribution, and be sure to specifically note its center.
  Make sure to include a plot of the distribution in your answer.
\end{enumerate}

\begin{Shaded}
\begin{Highlighting}[]
\FunctionTok{nrow}\NormalTok{(sample\_props50)}
\end{Highlighting}
\end{Shaded}

\begin{verbatim}
## [1] 15000
\end{verbatim}

\begin{Shaded}
\begin{Highlighting}[]
\NormalTok{sample\_props50 }\SpecialCharTok{\%\textgreater{}\%}
  \FunctionTok{ungroup}\NormalTok{() }\SpecialCharTok{\%\textgreater{}\%}
  \FunctionTok{summarise}\NormalTok{(}\AttributeTok{mean =} \FunctionTok{mean}\NormalTok{(p\_hat), }\AttributeTok{sd =} \FunctionTok{sd}\NormalTok{(p\_hat), }\AttributeTok{min =} \FunctionTok{min}\NormalTok{(p\_hat), }\AttributeTok{max =} \FunctionTok{max}\NormalTok{(p\_hat))}
\end{Highlighting}
\end{Shaded}

\begin{verbatim}
## # A tibble: 1 x 4
##    mean     sd   min   max
##   <dbl>  <dbl> <dbl> <dbl>
## 1 0.199 0.0568  0.02  0.42
\end{verbatim}

\begin{Shaded}
\begin{Highlighting}[]
\FunctionTok{ggplot}\NormalTok{(}\AttributeTok{data =}\NormalTok{ sample\_props50, }\FunctionTok{aes}\NormalTok{(}\AttributeTok{x =}\NormalTok{ p\_hat)) }\SpecialCharTok{+}
  \FunctionTok{geom\_histogram}\NormalTok{(}\AttributeTok{binwidth =} \FloatTok{0.02}\NormalTok{) }\SpecialCharTok{+}
  \FunctionTok{labs}\NormalTok{(}
    \AttributeTok{x =} \StringTok{"p\_hat (Doesn\textquotesingle{}t benefit)"}\NormalTok{,}
    \AttributeTok{title =} \StringTok{"Sampling distribution of p\_hat"}\NormalTok{,}
    \AttributeTok{subtitle =} \StringTok{"Sample size = 50, Number of samples = 15000"}
\NormalTok{  )}
\end{Highlighting}
\end{Shaded}

\pandocbounded{\includegraphics[keepaspectratio]{Lab5a_sampling_distributions_files/figure-latex/unnamed-chunk-9-1.pdf}}

\textbf{Insert your answer here}

\texttt{sample\_props50} has 15000 elements. The data can be visually
understood a being distributed normally centered around 0.2.

\subsection{Interlude: Sampling
distributions}\label{interlude-sampling-distributions}

The idea behind the \texttt{rep\_sample\_n} function is
\emph{repetition}. Earlier, you took a single sample of size \texttt{n}
(50) from the population of all people in the population. With this new
function, you can repeat this sampling procedure \texttt{rep} times in
order to build a distribution of a series of sample statistics, which is
called the \textbf{sampling distribution}.

Note that in practice one rarely gets to build true sampling
distributions, because one rarely has access to data from the entire
population.

Without the \texttt{rep\_sample\_n} function, this would be painful. We
would have to manually run the following code 15,000 times

\begin{Shaded}
\begin{Highlighting}[]
\NormalTok{global\_monitor }\SpecialCharTok{\%\textgreater{}\%}
  \FunctionTok{sample\_n}\NormalTok{(}\AttributeTok{size =} \DecValTok{50}\NormalTok{, }\AttributeTok{replace =} \ConstantTok{TRUE}\NormalTok{) }\SpecialCharTok{\%\textgreater{}\%}
  \FunctionTok{count}\NormalTok{(scientist\_work) }\SpecialCharTok{\%\textgreater{}\%}
  \FunctionTok{mutate}\NormalTok{(}\AttributeTok{p\_hat =}\NormalTok{ n }\SpecialCharTok{/}\FunctionTok{sum}\NormalTok{(n)) }\SpecialCharTok{\%\textgreater{}\%}
  \FunctionTok{filter}\NormalTok{(scientist\_work }\SpecialCharTok{==} \StringTok{"Doesn\textquotesingle{}t benefit"}\NormalTok{)}
\end{Highlighting}
\end{Shaded}

\begin{verbatim}
## # A tibble: 1 x 3
##   scientist_work      n p_hat
##   <chr>           <int> <dbl>
## 1 Doesn't benefit    11  0.22
\end{verbatim}

as well as store the resulting sample proportions each time in a
separate vector.

Note that for each of the 15,000 times we computed a proportion, we did
so from a \textbf{different} sample!

\begin{enumerate}
\def\labelenumi{\arabic{enumi}.}
\setcounter{enumi}{4}
\tightlist
\item
  To make sure you understand how sampling distributions are built, and
  exactly what the \texttt{rep\_sample\_n} function does, try modifying
  the code to create a sampling distribution of \textbf{25 sample
  proportions} from \textbf{samples of size 10}, and put them in a data
  frame named \texttt{sample\_props\_small}. Print the output. How many
  observations are there in this object called
  \texttt{sample\_props\_small}? What does each observation represent?
\end{enumerate}

\textbf{Insert your answer here}

\begin{Shaded}
\begin{Highlighting}[]
\NormalTok{sample\_props\_small }\OtherTok{\textless{}{-}}\NormalTok{ global\_monitor }\SpecialCharTok{\%\textgreater{}\%}
  \FunctionTok{rep\_sample\_n}\NormalTok{(}\AttributeTok{size =} \DecValTok{10}\NormalTok{, }\AttributeTok{reps =} \DecValTok{25}\NormalTok{, }\AttributeTok{replace =} \ConstantTok{TRUE}\NormalTok{) }\SpecialCharTok{\%\textgreater{}\%}
  \FunctionTok{count}\NormalTok{(scientist\_work) }\SpecialCharTok{\%\textgreater{}\%}
  \FunctionTok{mutate}\NormalTok{(}\AttributeTok{p\_hat =}\NormalTok{ n }\SpecialCharTok{/}\FunctionTok{sum}\NormalTok{(n)) }\SpecialCharTok{\%\textgreater{}\%}
  \FunctionTok{filter}\NormalTok{(scientist\_work }\SpecialCharTok{==} \StringTok{"Doesn\textquotesingle{}t benefit"}\NormalTok{)}

\NormalTok{sample\_props\_small}
\end{Highlighting}
\end{Shaded}

\begin{verbatim}
## # A tibble: 22 x 4
## # Groups:   replicate [22]
##    replicate scientist_work      n p_hat
##        <int> <chr>           <int> <dbl>
##  1         1 Doesn't benefit     2   0.2
##  2         2 Doesn't benefit     4   0.4
##  3         3 Doesn't benefit     5   0.5
##  4         5 Doesn't benefit     2   0.2
##  5         6 Doesn't benefit     1   0.1
##  6         7 Doesn't benefit     5   0.5
##  7         9 Doesn't benefit     1   0.1
##  8        10 Doesn't benefit     1   0.1
##  9        11 Doesn't benefit     2   0.2
## 10        12 Doesn't benefit     4   0.4
## # i 12 more rows
\end{verbatim}

\texttt{sample\_props\_small} has 20 observations. Each one is the
sample proportion of ``Doesn't benefit'' (\texttt{p\_hat}) from one
random sample of 10 people.

\subsection{Sample size and the sampling
distribution}\label{sample-size-and-the-sampling-distribution}

Mechanics aside, let's return to the reason we used the
\texttt{rep\_sample\_n} function: to compute a sampling distribution,
specifically, the sampling distribution of the proportions from samples
of 50 people.

\begin{Shaded}
\begin{Highlighting}[]
\FunctionTok{ggplot}\NormalTok{(}\AttributeTok{data =}\NormalTok{ sample\_props50, }\FunctionTok{aes}\NormalTok{(}\AttributeTok{x =}\NormalTok{ p\_hat)) }\SpecialCharTok{+}
  \FunctionTok{geom\_histogram}\NormalTok{(}\AttributeTok{binwidth =} \FloatTok{0.02}\NormalTok{)}
\end{Highlighting}
\end{Shaded}

The sampling distribution that you computed tells you much about
estimating the true proportion of people who think that the work
scientists do doesn't benefit them. Because the sample proportion is an
unbiased estimator, the sampling distribution is centered at the true
population proportion, and the spread of the distribution indicates how
much variability is incurred by sampling only 50 people at a time from
the population.

In the remainder of this section, you will work on getting a sense of
the effect that sample size has on your sampling distribution.

\begin{enumerate}
\def\labelenumi{\arabic{enumi}.}
\setcounter{enumi}{5}
\tightlist
\item
  Use the app below to create sampling distributions of proportions of
  \emph{Doesn't benefit} from samples of size 10, 50, and 100. Use 5,000
  simulations. What does each observation in the sampling distribution
  represent? How does the mean, standard error, and shape of the
  sampling distribution change as the sample size increases? How (if at
  all) do these values change if you increase the number of simulations?
  (You do not need to include plots in your answer.)
\end{enumerate}

\textbf{Insert your answer here}

The observations in the sampling distribution are the proportions of
responses ``Doesn't benefit'' in each randomly sampled set of size n.~As
the sample size is increased from 10 to 50 to 100, the mean value
remains fairly constant, and close to the true proportion of 0.20. The
mean value for n=10 is somewhat higher than 0.20 since any sample
containing no response ``Doesn't benefit'' is rejected during the
simulation process. The standard deviation decreases and follows very
closely the theoretical value of √(p(1-p)/n). The distribution is first
lumpy and right-skewed (for n=10), then bell-shaped and approximately
normally distributed for n=50 and n=100. Increase in the number of
simulations has no effect on the mean and standard deviation, only makes
the histogram smoother since the spread of the sampling distribution
depends only on the sample size.

\begin{center}\rule{0.5\linewidth}{0.5pt}\end{center}

\subsection{More Practice}\label{more-practice}

So far, you have only focused on estimating the proportion of those you
think the work scientists doesn't benefit them. Now, you'll try to
estimate the proportion of those who think it does.

Note that while you might be able to answer some of these questions
using the app, you are expected to write the required code and produce
the necessary plots and summary statistics. You are welcome to use the
app for exploration.

\begin{enumerate}
\def\labelenumi{\arabic{enumi}.}
\setcounter{enumi}{6}
\tightlist
\item
  Take a sample of size 15 from the population and calculate the
  proportion of people in this sample who think the work scientists do
  enhances their lives. Using this sample, what is your best point
  estimate of the population proportion of people who think the work
  scientists do enchances their lives?
\end{enumerate}

\textbf{Insert your answer here}

\begin{Shaded}
\begin{Highlighting}[]
\NormalTok{samp15 }\OtherTok{\textless{}{-}}\NormalTok{ global\_monitor }\SpecialCharTok{\%\textgreater{}\%}
  \FunctionTok{sample\_n}\NormalTok{(}\DecValTok{15}\NormalTok{)}

\NormalTok{samp15 }\SpecialCharTok{\%\textgreater{}\%}
  \FunctionTok{count}\NormalTok{(scientist\_work) }\SpecialCharTok{\%\textgreater{}\%}
  \FunctionTok{mutate}\NormalTok{(}\AttributeTok{p =}\NormalTok{ n }\SpecialCharTok{/} \FunctionTok{sum}\NormalTok{(n))}
\end{Highlighting}
\end{Shaded}

\begin{verbatim}
## # A tibble: 1 x 3
##   scientist_work     n     p
##   <chr>          <int> <dbl>
## 1 Benefits          15     1
\end{verbatim}

\begin{Shaded}
\begin{Highlighting}[]
\NormalTok{p\_15 }\OtherTok{\textless{}{-}} \FunctionTok{mean}\NormalTok{(samp15}\SpecialCharTok{$}\NormalTok{scientist\_work }\SpecialCharTok{==} \StringTok{"Benefits"}\NormalTok{)}
\end{Highlighting}
\end{Shaded}

In this sample of 15, the proportion who think the work scientists do
benefits them is 0.8. The sample proportion is the best point estimate
of the population proportion, so my estimate is 0.8.

\begin{enumerate}
\def\labelenumi{\arabic{enumi}.}
\setcounter{enumi}{7}
\tightlist
\item
  Since you have access to the population, simulate the sampling
  distribution of proportion of those who think the work scientists do
  enchances their lives for samples of size 15 by taking 2000 samples
  from the population of size 15 and computing 2000 sample proportions.
  Store these proportions in as \texttt{sample\_props15}. Plot the data,
  then describe the shape of this sampling distribution. Based on this
  sampling distribution, what would you guess the true proportion of
  those who think the work scientists do enchances their lives to be?
  Finally, calculate and report the population proportion.
\end{enumerate}

\textbf{Insert your answer here}

\begin{Shaded}
\begin{Highlighting}[]
\NormalTok{sample\_props15 }\OtherTok{\textless{}{-}}\NormalTok{ global\_monitor }\SpecialCharTok{\%\textgreater{}\%}
  \FunctionTok{rep\_sample\_n}\NormalTok{(}\AttributeTok{size =} \DecValTok{15}\NormalTok{, }\AttributeTok{reps =} \DecValTok{2000}\NormalTok{, }\AttributeTok{replace =} \ConstantTok{TRUE}\NormalTok{) }\SpecialCharTok{\%\textgreater{}\%}
  \FunctionTok{count}\NormalTok{(scientist\_work) }\SpecialCharTok{\%\textgreater{}\%}
  \FunctionTok{mutate}\NormalTok{(}\AttributeTok{p\_hat =}\NormalTok{ n }\SpecialCharTok{/} \FunctionTok{sum}\NormalTok{(n)) }\SpecialCharTok{\%\textgreater{}\%}
  \FunctionTok{filter}\NormalTok{(scientist\_work }\SpecialCharTok{==} \StringTok{"Benefits"}\NormalTok{)}
\end{Highlighting}
\end{Shaded}

\begin{Shaded}
\begin{Highlighting}[]
\FunctionTok{ggplot}\NormalTok{(sample\_props15, }\FunctionTok{aes}\NormalTok{(}\AttributeTok{x =}\NormalTok{ p\_hat)) }\SpecialCharTok{+}
  \FunctionTok{geom\_histogram}\NormalTok{(}\AttributeTok{binwidth =} \DecValTok{1}\SpecialCharTok{/}\DecValTok{15}\NormalTok{, }\AttributeTok{center =} \DecValTok{0}\NormalTok{)}
\end{Highlighting}
\end{Shaded}

\pandocbounded{\includegraphics[keepaspectratio]{Lab5a_sampling_distributions_files/figure-latex/unnamed-chunk-14-1.pdf}}

\begin{Shaded}
\begin{Highlighting}[]
\FunctionTok{mean}\NormalTok{(sample\_props15}\SpecialCharTok{$}\NormalTok{p\_hat)}
\end{Highlighting}
\end{Shaded}

\begin{verbatim}
## [1] 0.8038
\end{verbatim}

\begin{Shaded}
\begin{Highlighting}[]
\FunctionTok{sd}\NormalTok{(sample\_props15}\SpecialCharTok{$}\NormalTok{p\_hat)}
\end{Highlighting}
\end{Shaded}

\begin{verbatim}
## [1] 0.1046259
\end{verbatim}

\begin{Shaded}
\begin{Highlighting}[]
\NormalTok{global\_monitor }\SpecialCharTok{\%\textgreater{}\%}
  \FunctionTok{count}\NormalTok{(scientist\_work) }\SpecialCharTok{\%\textgreater{}\%}
  \FunctionTok{mutate}\NormalTok{(}\AttributeTok{p =}\NormalTok{ n }\SpecialCharTok{/} \FunctionTok{sum}\NormalTok{(n))}
\end{Highlighting}
\end{Shaded}

\begin{verbatim}
## # A tibble: 2 x 3
##   scientist_work      n     p
##   <chr>           <int> <dbl>
## 1 Benefits        80000   0.8
## 2 Doesn't benefit 20000   0.2
\end{verbatim}

The shape of the distribution will be unimodal and left-skewed because
it will have a longer tail at the low proportions end. It does not fit
the normal distribution since the true proportion (0.80) is close to 1,
and the sample size is also small (n = 15, thus n(1 - p) = 3 \textless{}
10). It will be around 0.804, so my guess will be 0.80. The actual
proportion is 0.8.

\begin{enumerate}
\def\labelenumi{\arabic{enumi}.}
\setcounter{enumi}{8}
\tightlist
\item
  Change your sample size from 15 to 150, then compute the sampling
  distribution using the same method as above, and store these
  proportions in a new object called \texttt{sample\_props150}. Describe
  the shape of this sampling distribution and compare it to the sampling
  distribution for a sample size of 15. Based on this sampling
  distribution, what would you guess to be the true proportion of those
  who think the work scientists do enchances their lives?
\end{enumerate}

\textbf{Insert your answer here}

\begin{Shaded}
\begin{Highlighting}[]
\NormalTok{sample\_props150 }\OtherTok{\textless{}{-}}\NormalTok{ global\_monitor }\SpecialCharTok{\%\textgreater{}\%}
  \FunctionTok{rep\_sample\_n}\NormalTok{(}\AttributeTok{size =} \DecValTok{150}\NormalTok{, }\AttributeTok{reps =} \DecValTok{2000}\NormalTok{, }\AttributeTok{replace =} \ConstantTok{TRUE}\NormalTok{) }\SpecialCharTok{\%\textgreater{}\%}
  \FunctionTok{count}\NormalTok{(scientist\_work) }\SpecialCharTok{\%\textgreater{}\%}
  \FunctionTok{mutate}\NormalTok{(}\AttributeTok{p\_hat =}\NormalTok{ n }\SpecialCharTok{/} \FunctionTok{sum}\NormalTok{(n)) }\SpecialCharTok{\%\textgreater{}\%}
  \FunctionTok{filter}\NormalTok{(scientist\_work }\SpecialCharTok{==} \StringTok{"Benefits"}\NormalTok{)}

\FunctionTok{ggplot}\NormalTok{(sample\_props150, }\FunctionTok{aes}\NormalTok{(}\AttributeTok{x =}\NormalTok{ p\_hat)) }\SpecialCharTok{+}
  \FunctionTok{geom\_histogram}\NormalTok{(}\AttributeTok{binwidth =} \FloatTok{0.02}\NormalTok{)}
\end{Highlighting}
\end{Shaded}

\pandocbounded{\includegraphics[keepaspectratio]{Lab5a_sampling_distributions_files/figure-latex/unnamed-chunk-17-1.pdf}}

\begin{Shaded}
\begin{Highlighting}[]
\FunctionTok{mean}\NormalTok{(sample\_props150}\SpecialCharTok{$}\NormalTok{p\_hat)}
\end{Highlighting}
\end{Shaded}

\begin{verbatim}
## [1] 0.8000933
\end{verbatim}

\begin{Shaded}
\begin{Highlighting}[]
\FunctionTok{sd}\NormalTok{(sample\_props150}\SpecialCharTok{$}\NormalTok{p\_hat)}
\end{Highlighting}
\end{Shaded}

\begin{verbatim}
## [1] 0.03282338
\end{verbatim}

The sampling distribution of n = 150 is symmetric and bell-shaped
(normally distributed) as opposed to being left skewed in the case of n
= 15. The mean of the two distributions is the same, but the standard
deviation is smaller for n = 150. My estimate of the population
proportion is 0.80.

\begin{enumerate}
\def\labelenumi{\arabic{enumi}.}
\setcounter{enumi}{9}
\tightlist
\item
  Of the sampling distributions from 2 and 3, which has a smaller
  spread? If you're concerned with making estimates that are more often
  close to the true value, would you prefer a sampling distribution with
  a large or small spread?\\
  \strut \\
  The n = 150 distribution has the smaller spread (standard deviation
  0.033 versus 0.103). To make estimates that are more often close to
  the true value, I'd prefer a sampling distribution with a small
  spread, because then a typical sample proportion lands close to the
  population proportion.
\end{enumerate}

\begin{center}\rule{0.5\linewidth}{0.5pt}\end{center}

\end{document}
